\section{Design principles}\label{sec:principles}

Every decision in the benchmark traces to one of seven principles.

\begin{description}[style=unboxed,leftmargin=0pt,itemsep=2pt]
\item[Luck-robust.] A score that does not deflate for the number of strategies tried is a score of the best overfitter, so the deflated Sharpe is a gate and each agent's declared trials raise its own bar.
\item[Reliability-gated.] An agent that is safe on average is not safe, so eligibility requires clearing the per-run bar on every seed and every window, which is a stronger claim than having an edge and is meant to be.
\item[Process-gated.] A return earned by bypassing a risk control does not count, and one block-severity violation in any run zeroes the entry.
\item[Point-in-time.] The agent interface can express only history at or before the decision, and that boundary is the thing to audit.
\item[Deterministic.] A score is a pure function of a field of submissions and a configuration, byte-identical on the tested platforms, and the continuous-integration matrix checks it.
\item[Forward-attested.] A commitment binds an artifact before an operator-declared deadline, and a chain lets anyone verify the resulting byte history. The forward interpretation still trusts the operator's wall-time mapping, held-out-data custody and claim of non-observation before commitment. If the host enters a future board, its entries are marked and governed by the same declared cap; no model-entry board exists yet.
\item[Judge-free.] No model grades another, and every gate is an assertion over numbers, which is what makes the self-audit executable.
\end{description}

\Cref{tab:design-constraints} restates each principle as the requirement the artifact enforces and the failure a violation would permit.

\begin{table}[t]
\centering
\footnotesize
\setlength{\tabcolsep}{4pt}
\caption{Design constraints. Each principle is stated as the requirement the kernel or protocol enforces and the failure mode a benchmark that violated it would exhibit. The requirement column names mechanisms; whether each is exercised by the frozen evidence is stated in the section that describes it.}
\label{tab:design-constraints}
\begin{tabular}{>{\raggedright\arraybackslash}p{0.15\linewidth}>{\raggedright\arraybackslash}p{0.4\linewidth}>{\raggedright\arraybackslash}p{0.38\linewidth}}
\toprule
Principle & Requirement & Failure mode if violated \\
\midrule
Luck-robust & DSR gate at an effective trial count no smaller than the observed field, raised by each agent's declared trials (\cref{sec:stats}) & the best overfitter of a large search ranks first \\
Reliability-gated & every run, on every seed and every window, clears the per-run PSR bar under the configured verdict (\cref{sec:predicate}) & an agent that is safe on average or in one regime ranks as safe \\
Process-gated & one block-severity event in any run makes the entry ineligible (\cref{sec:predicate}) & a return earned by bypassing a risk control buys rank \\
Point-in-time & the observation holds only bars at or before the decision index (\cref{sec:integrity}) & look-ahead inflates returns without any visible defect \\
Deterministic & the score is a pure function of field and configuration, with golden fixtures checked on three operating systems (\cref{sec:integrity}) & two operators disagree about one submission and tampering is undetectable \\
Forward-attested & a salted commitment to the artifact digest before a declared deadline, on a chained signed board (\cref{sec:attest}) & a strategy tuned after its evaluation window was seen is scored as if it had not been \\
Judge-free & every gate is an assertion over numbers (\cref{sec:selfaudit}) & scores inherit a grader's biases and the self-audit cannot be executed \\
\bottomrule
\end{tabular}
\end{table}

Two consequences bear on the experiments. First, what determinism buys and what it does not. A deterministic kernel supports replicability and auditability: anyone can recompute every score, and tampering within the declared digest coverage of \cref{sec:integrity} is detectable by a verifier holding the manifest. It guarantees nothing about validity: the first finding shows the benchmark shipping with a deflation prior off by up to a factor of $\sqrt{8760}$ while the eight-attack self-audit stayed green, because a reproducible scorer is wrong in exactly the same way everywhere. The self-audit tests gate evasion, not gate calibration. A PSR near one beside a DSR near zero is mathematically possible because the statistics test different thresholds; when it appeared we inspected the implied annualized DSR threshold, whose dimensional scale revealed the error. \Cref{sec:units} states that calibration check explicitly. Second, the reliability and process principles together mean the benchmark can decline to certify any agent at all on a dataset, including the market itself. Whether universal refusal is the right outcome on a given history is a question \cref{sec:experiments} puts to the data.

\paragraph{What validates the instrument, and what does not.} There is no prior trading-agent benchmark for SharpeBench to agree with, and no human or practitioner rating study has been run (\cref{sec:limitations}), so the construct-validity evidence here is discriminative: entrants whose properties are known by construction must land where the construct says they should. Four checks of that kind are reported. Seeded zero-skill agents never beat a reference agent (\cref{sec:falsify,sec:luck1000}); a synthetic agent with an injected edge becomes eligible above a sampled boundary and not below it (\cref{sec:witness}); ten constructed attacks are demoted on every commit (\cref{sec:selfaudit}); and a deliberately fragile agent separates from buy-and-hold under perturbation in a unit test (\cref{sec:perturb}). None of them validates the gates' calibration on real agents. The witness is one common-random-number draw; the committed perturbation report on a real panel separates nothing; no real or third-party agent has been admitted; the refusal has low power at the shipped window geometry (\cref{sec:power}); and the self-audit tests evasion rather than calibration, as the first finding shows. The protocol is therefore validated as a procedure that behaves correctly on known entrants, not as a measurement whose verdicts on unknown agents have been checked against anything external.
